\section{Jones-Formalism}

The Jones formalism can be used to describe mathematically the polarization state of light and the effect of polarizers and wave plates.
Here the polarization state is represented as a two-dimensional vector and the optical elements are represented as a 2x2 matrix. 
\begin{equation*}
\begin{aligned}
\begin{array}{lcl}
\text{Linearly polarized light in x-direction:} &
\vec{J}_x = &
\begin{pmatrix} 1 \\ 0 \end{pmatrix}
\\[0.5em]

\text{Linearly polarized light in y-direction:} &
\vec{J}_y = &
\begin{pmatrix} 0 \\ 1 \end{pmatrix}
\\[0.5em]

\text{Linearly polarized light at }45^\circ\text{:} &
\vec{J}_{45} = &
\frac{1}{\sqrt{2}}
\begin{pmatrix} 1 \\ 1 \end{pmatrix}
\\[0.5em]

\text{Right circularly polarized light:} &
\vec{J}_R = &
\frac{1}{\sqrt{2}}
\begin{pmatrix} 1 \\ -i \end{pmatrix}
\\[0.5em]

\text{Left circularly polarized light:} &
\vec{J}_L = &
\frac{1}{\sqrt{2}}
\begin{pmatrix} 1 \\ i \end{pmatrix}
\end{array}
\end{aligned}
\end{equation*}

The Jones matrix for the $\lambda/4$-plate with the fast axis along the x-direction is given by:
\begin{equation*}
\mathbf{M}_{\lambda/4} = \frac{1}{\sqrt{2}}
\begin{pmatrix} 1-i & 0 \\
 0 & 1 + i \end{pmatrix}
\end{equation*}
This matrix introduces a phase shift of $\pi/2$ between the x and y components of the electric field, which is characteristic of a quarter-wave plate.
When linearly polarized light at 45 degrees passes through this $\lambda/4$-plate, it becomes circularly polarized.
Conversely, circularly polarized light can be converted back into linearly polarized light by passing it through the same $\lambda/4$-plate
swith the appropriate orientation.